\chapter{Metodología y planificación del trabajo}
\label{chap:metodologia}

\lettrine{E}{l} desarrollo de sistemas \acrshort{bci} integrados presenta desafíos únicos que requieren enfoques metodológicos adaptados a la naturaleza experimental del dominio. La complejidad inherente de integrar captura de señales biomédicas con eventos de entrada del usuario, junto con los requisitos de tiempo real, demanda una metodología que combine rigor en la planificación con flexibilidad para adaptarse a descubrimientos durante el desarrollo.

Este capítulo presenta la metodología empleada, comenzando por la planificación inicial del proyecto incluyendo estimaciones temporales y económicas, seguido del seguimiento realizado durante el desarrollo, y finalizando con el análisis de desvíos respecto a la planificación original. La documentación de estos aspectos proporciona perspectivas valiosas sobre la gestión de proyectos de investigación en el dominio \acrshort{bci}.

\section{Enfoque metodológico}
\label{sec:enfoque_metodologico}

Este proyecto se ha desarrollado siguiendo un modelo \textbf{incremental-iterativo} que combina las ventajas de una definición temprana de requisitos con la flexibilidad necesaria para refinar el resultado en ciclos sucesivos. Esta metodología emergió como la más apropiada para abordar la complejidad inherente del desarrollo de sistemas \acrshort{bci}, donde la experimentación y validación temprana son fundamentales para el éxito del proyecto.

La \textbf{incrementalidad} permite que el \textit{software} se construya por módulos funcionales independientes, cada uno de los cuales añade nuevas prestaciones al sistema global sin comprometer la estabilidad de componentes previamente desarrollados. Esta aproximación ha sido especialmente valiosa en un proyecto donde diferentes subsistemas pueden desarrollarse en paralelo y validarse de forma independiente.

La \textbf{iteración} garantiza que cada módulo se somete a múltiples rondas de desarrollo, validación y refactorización, permitiendo la detección rápida de errores y su corrección antes de que se propaguen a otros componentes del sistema. Esta característica resultó fundamental durante las fases de integración, donde se identificaron limitaciones arquitectónicas que requirieron refactorización significativa.

Aunque la metodología base se mantuvo como incremental-iterativa a lo largo del proyecto, la estrategia de implementación se adaptó para responder a desafíos específicos como la migración temprana de Julia a Python, los múltiples intentos fallidos de integración directa entre módulos, y la eventual adopción de arquitecturas \acrshort{api} como solución robusta para la comunicación entre componentes.

La validación experimental continua mediante pruebas manuales iterativas ha proporcionado retroalimentación constante sobre el funcionamiento del sistema. En cada iteración se realizaba validación manual de funcionalidades, verificando que el comportamiento correspondía con las expectativas antes de proceder con nuevas funcionalidades. Esta validación incluía pruebas con dispositivos sintéticos, controladores físicos y cascos \acrshort{eeg} reales para garantizar que las decisiones de diseño estuvieran fundamentadas en evidencia práctica.

\subsection{Estructura de iteraciones}
\label{subsec:estructura_iteraciones}

Cada iteración sigue un ciclo de desarrollo que incluye las siguientes fases:

\begin{itemize}
    \item \textbf{Análisis de requisitos}: Identificación y refinamiento de funcionalidades necesarias basado en el conocimiento adquirido y las limitaciones descubiertas en iteraciones previas.
    \item \textbf{Diseño}: Definición de la arquitectura de componentes y las interfaces necesarias para implementar las funcionalidades identificadas.
    \item \textbf{Implementación}: Desarrollo de código siguiendo principios \acrshort{solid}~\cite{martin2003agile}, con énfasis en la modularidad y extensibilidad.
    \item \textbf{Pruebas manuales}: Validación exhaustiva de funcionalidades mediante pruebas manuales con dispositivos reales y sintéticos, verificando tanto el funcionamiento técnico como la correspondencia con el comportamiento esperado.
    \item \textbf{Evaluación y planificación}: Análisis de resultados obtenidos, identificación de problemas o limitaciones, y planificación de la siguiente iteración o módulo a desarrollar.
\end{itemize}

Este enfoque pragmático permitió una adaptación continua a los desafíos encontrados, priorizando la funcionalidad práctica sobre protocolos formales de \textit{testing} automatizado como \glspl{unit_test}, \glspl{integration_test} o metodologías como \gls{tdd}, lo cual resultó apropiado para la naturaleza experimental del proyecto.

\section{Planificación inicial del proyecto}
\label{sec:planificacion_inicial}

La planificación inicial del proyecto contemplaba un desarrollo más lineal y predecible, basado en la asunción de que la integración directa entre módulos desarrollados independientemente sería relativamente directa. Esta planificación inicial subestimó la complejidad de los desafíos de integración en sistemas \acrshort{bci}.

\subsection{Fases principales planificadas}
\label{subsec:fases_planificadas}

El proyecto se planificó inicialmente en 5 fases principales distribuidas a lo largo de 16 semanas, siguiendo la metodología incremental-iterativa establecida. Las fases principales incluían: documentación e investigación, análisis y diseño del sistema, implementación del módulo de captura de \textit{inputs}, implementación del módulo de adquisición \acrshort{eeg}, y desarrollo del sistema integrado de registro y sincronización.

\begin{enumerate}
    \item \textbf{Documentación e investigación}: Estudio del estado del arte en sistemas \acrshort{bci}, procesamiento de señales \acrshort{eeg} y análisis de bibliotecas disponibles.
    \item \textbf{Análisis y diseño del sistema}: Definición de la arquitectura, interfaces entre módulos y estructuras de datos.
    \item \textbf{Implementación del módulo de captura de \textit{inputs}}: Desarrollo del sistema para registrar eventos de teclado, ratón y controladores.
    \item \textbf{Implementación del módulo de adquisición \acrshort{eeg}}: Integración con BrainFlow y configuración para diversos dispositivos.
    \item \textbf{Desarrollo del sistema integrado}: Implementación del componente unificado para capturar y alinear temporalmente las señales \acrshort{eeg} con eventos de entrada, incluyendo almacenamiento estructurado de datos sincronizados en múltiples formatos estándar.
\end{enumerate}

\subsection{Cronograma inicial planificado}
\label{subsec:cronograma_inicial}

La planificación temporal inicial establecía objetivos específicos y entregables para cada iteración, organizados en un desarrollo incremental donde cada fase construye sobre las anteriores.

\begin{table}[htbp!]
  \centering
  \resizebox{\textwidth}{!}{
    \rowcolors{2}{white}{udcgray!25}
    \begin{tabular}{c|l|l|c}
    \rowcolor{udcpink!25}
    \textbf{Iteración} & \textbf{Objetivo principal} & \textbf{Entregables} & \textbf{Duración} \\
    \hline
    Iter. 0 & Investigación y análisis & Estado del arte, especificación inicial & 4 semanas \\
    Iter. 1 & Módulo captura \textit{inputs} & Prototipo funcional + \acrshort{gui} & 2 semanas \\
    Iter. 2 & Módulo procesamiento \acrshort{eeg} & Captura básica + visualización & 3 semanas \\
    Iter. 3 & Integración mediante \acrshort{lsl} & Sistema unificado funcional & 3 semanas \\
    Iter. 4 & Funcionalidades avanzadas & Análisis avanzado + exportación & 3 semanas \\
    Iter. 5 & Finalización del sistema & Sistema completo y memoria iniciada & 1 semana \\
    \end{tabular}
  }
  \caption{Planificación inicial de iteraciones del proyecto}
  \label{tab:planificacion_iteraciones}
\end{table}

La Tabla~\ref{tab:planificacion_iteraciones} muestra la distribución temporal planificada, donde cada iteración incluía las fases típicas de análisis, diseño, implementación, pruebas y evaluación. La documentación se planificó como una actividad transversal a todo el proyecto, con una fase final para finalizar el sistema y comenzar la elaboración de la memoria del TFG.

\subsection{Diagrama de Gantt inicial}
\label{subsec:gantt_inicial}

La planificación inicial asumía que la fase de integración directa mediante \acrshort{lsl} sería relativamente breve y no problemática, asignando únicamente 3 semanas para esta actividad crítica. Esta subestimación refleja la dificultad inherente de predecir la complejidad de integración en sistemas experimentales.

\begin{sidewaysfigure}[p]
\centering
\resizebox{0.9\textwidth}{0.7\textheight}{
\begin{ganttchart}[
    canvas/.append style={fill=none, draw=black!5, line width=.75pt},
    hgrid style/.style={draw=black!5, line width=.75pt},
    vgrid={*1{draw=black!5, line width=.75pt}},
    title/.style={draw=none, fill=none},
    title label font=\bfseries\small,
    title label node/.append style={below=5pt},
    include title in canvas=false,
    bar label font=\mdseries\small\color{black!70},
    bar label node/.append style={left=2cm},
    bar/.append style={draw=none, fill=black!63},
    milestone/.append style={fill=yellow, draw=black, diamond},
    bar progress label font=\mdseries\small\color{black!70},
    group incomplete/.append style={fill=groupblue},
    group left shift=0,
    group right shift=0,
    group height=.5,
    group peaks tip position=0,
    group label node/.append style={left=1cm},
    group progress label font=\bfseries\small
  ]{1}{64}
  
  \gantttitle{Marzo 2025}{16} \gantttitle{Abril 2025}{16} \gantttitle{Mayo 2025}{16} \gantttitle{Junio 2025}{16} \\
  \gantttitlelist{1,8,15,22}{4} \gantttitlelist{29,5,12,19}{4} \gantttitlelist{26,3,10,17}{4} \gantttitlelist{24,31,7,14}{4} \\
  
  \ganttgroup{Iter. 0: Investigación}{1}{16} \\
  \ganttbar{Análisis de requisitos}{1}{4} \\
  \ganttbar{Estado del arte}{5}{12} \\
  \ganttbar{Diseño arquitectónico}{13}{16} \\
  \ganttmilestone{Especificación completa}{16} \\
  
  \ganttgroup{Iter. 1: Módulo \textit{inputs}}{17}{24} \\
  \ganttbar{Análisis técnico}{17}{18} \\
  \ganttbar{Implementación núcleo}{19}{22} \\
  \ganttbar{Desarrollo \acrshort{gui}}{21}{24} \\
  \ganttbar{Pruebas integración}{23}{24} \\
  \ganttmilestone{Módulo \textit{inputs} completo}{24} \\
  
  \ganttgroup{Iter. 2: Módulo \acrshort{eeg}}{25}{36} \\
  \ganttbar{Diseño arquitectura}{25}{26} \\
  \ganttbar{Implementación captura}{27}{32} \\
  \ganttbar{Visualización tiempo real}{31}{34} \\
  \ganttbar{Pruebas con dispositivos}{35}{36} \\
  \ganttmilestone{Módulo \acrshort{eeg} funcional}{36} \\
  
  \ganttgroup{Iter. 3: Integración \acrshort{lsl}}{37}{48} \\
  \ganttbar{Análisis integración}{37}{38} \\
  \ganttbar{Implementación sincronización}{39}{44} \\
  \ganttbar{Pruebas sistema completo}{45}{47} \\
  \ganttbar{Optimización rendimiento}{46}{48} \\
  \ganttmilestone{Sistema integrado}{48} \\
  
  \ganttgroup{Iter. 4: Funcionalidades avanzadas}{49}{60} \\
  \ganttbar{Algoritmos procesamiento}{49}{54} \\
  \ganttbar{Exportación \textit{datasets}}{53}{58} \\
  \ganttbar{Validación final}{59}{60} \\
  \ganttmilestone{Sistema completo}{60} \\
  
  \ganttgroup{Iter. 5: Documentación}{61}{64} \\
  \ganttbar{Memoria TFG}{61}{64} \\
  \ganttmilestone{Proyecto finalizado}{64} \\
  
\end{ganttchart}
}
\caption{Diagrama de Gantt - Planificación inicial del proyecto}
\label{fig:gantt_inicial}
\end{sidewaysfigure}

La Figura~\ref{fig:gantt_inicial} muestra la planificación inicial donde se desgrana cada fase en sus etapas típicas: análisis, diseño, implementación y pruebas, siguiendo las iteraciones de la metodología incremental-iterativa establecida.

\subsection{Planificación económica inicial}
\label{subsec:planificacion_economica}

La evaluación económica inicial del proyecto reflejaba una estrategia de optimización de costes basada en la utilización de recursos disponibles y herramientas de código abierto.

\subsubsection{Recursos humanos planificados}
\label{subsubsec:recursos_humanos_planificados}

Se presenta un supuesto de lo que costaría el desarrollo del sistema por un equipo profesional, basándose en salarios de mercado realistas para el contexto español. Esta estimación proporciona una valoración del esfuerzo y recursos necesarios para replicar el proyecto en un contexto comercial.

\begin{table}[htbp!]
  \centering
  \resizebox{\textwidth}{!}{
    \rowcolors{2}{white}{udcgray!25}
    \begin{tabular}{l|c|c|c|c}
    \rowcolor{udcpink!25}
    \textbf{Perfil} & \textbf{Horas planificadas} & \textbf{Coste/hora} & \textbf{Coste total} & \textbf{Descripción} \\
    \hline
    Desarrollador Junior & 280 horas & 12€/h & 3.360€ & Implementación y pruebas \\
    Analista de sistemas & 40 horas & 18€/h & 720€ & Análisis y diseño \\
    Especialista \acrshort{bci} & 20 horas & 20€/h & 400€ & Validación neurocientífica \\
    \hline
    \textbf{Total recursos humanos} & \textbf{340 horas} & & \textbf{4.480€} & \\
    \end{tabular}
  }
  \caption{Estimación de costes de recursos humanos}
  \label{tab:recursos_humanos_inicial}
\end{table}

La Tabla~\ref{tab:recursos_humanos_inicial} detalla los perfiles necesarios y sus costes asociados, calculados en base a salarios de mercado para el contexto español, asumiendo unas 40h de trabajo semanales.

\subsubsection{Costes totales estimados inicialmente}
\label{subsubsec:costes_totales_estimados}

El desglose completo de costes estimados para el proyecto incluye tanto los recursos humanos como los recursos materiales y tecnológicos necesarios, valorados a precios de mercado independientemente de su disponibilidad previa.

\begin{table}[htbp!]
  \centering
  \resizebox{\textwidth}{!}{
    \rowcolors{2}{white}{udcgray!25}
    \begin{tabular}{l|c|l}
    \rowcolor{udcpink!25}
    \textbf{Categoría} & \textbf{Coste estimado} & \textbf{Observaciones} \\
    \hline
    Recursos humanos & 4.480€ & Desarrollo completo del sistema \\
    \textit{Hardware} (amortización) & 800€ & PC desarrollo y casco \acrshort{eeg} (4 meses uso) \\
    \textit{Software} y herramientas & 0€ & Herramientas de código abierto \\
    Infraestructura & 0€ & Desarrollo local \\
    Contingencias (10\%) & 528€ & Imprevistos y sobrecostes \\
    \hline
    \textbf{Total proyecto} & \textbf{5.808€} & \\
    \end{tabular}
  }
  \caption{Estimación inicial de costes del proyecto}
  \label{tab:costes_inicial}
\end{table}

La Tabla~\ref{tab:costes_inicial} presenta la estimación económica inicial, incluyendo la amortización proporcional del \textit{hardware} utilizado durante el período de desarrollo. La valoración del \textit{hardware} incluye un PC de desarrollo (€1.500) y casco \acrshort{eeg} g.tec Unicorn (€1.500), amortizados durante los 4 meses de duración del proyecto.

\subsection{Análisis de riesgos inicial}
\label{subsec:riesgos_inicial}

La identificación proactiva de riesgos permitió desarrollar estrategias de mitigación anticipadas, aunque algunos riesgos no anticipados emergieron durante el desarrollo.

\begin{table}[htbp!]
  \centering
  \resizebox{\textwidth}{!}{
    \rowcolors{2}{white}{udcgray!25}
    \begin{tabular}{l|c|c|l}
    \rowcolor{udcpink!25}
    \textbf{Riesgo identificado} & \textbf{Prob.} & \textbf{Impacto} & \textbf{Estrategia de mitigación planificada} \\
    \hline
    Complejidad integración módulos & Media & Alta & Prototipado temprano con \acrshort{lsl} \\
    Rendimiento tiempo real & Media & Alta & Optimización mediante \gls{callback} \\
    Disponibilidad dispositivos \acrshort{eeg} & Baja & Alta & Dispositivos sintéticos BrainFlow \\
    Curva aprendizaje tecnologías & Media & Media & Investigación exhaustiva previa \\
    Dependencias bibliotecas externas & Baja & Media & Selección de bibliotecas estables \\
    \hline
    \end{tabular}
  }
  \caption{Análisis inicial de riesgos del proyecto}
  \label{tab:riesgos_inicial}
\end{table}

La Tabla~\ref{tab:riesgos_inicial} documenta los principales riesgos identificados en la planificación inicial junto con las estrategias de mitigación previstas.

\section{Seguimiento y control del proyecto}
\label{sec:seguimiento_control}

Durante el desarrollo se implementaron mecanismos de seguimiento que permitieron documentar el progreso real y detectar desvíos respecto a la planificación inicial. El control de versiones mediante Git proporcionó trazabilidad completa de los cambios, mientras que el sistema de \gls{pre-commit hooks} contribuyó a mantener la calidad del código. La implementación progresiva de \gls{type hints} facilitó el mantenimiento y la detección de errores.

\subsection{Hitos principales alcanzados}
\label{subsec:hitos_principales}

Los hitos del proyecto reflejaron tanto los éxitos como los desafíos encontrados durante el desarrollo:

\begin{itemize}
    \item \textbf{Hito 1 (19/04)}: Módulos básicos desarrollados, primeras limitaciones de integración identificadas
    \item \textbf{Hito 2 (30/04)}: Refactorización \acrshort{solid} completada, persistencia de problemas de integración
    \item \textbf{Hito 3 (22/05)}: \acrshort{api} de acceso a datos y explorador web implementados exitosamente, punto de inflexión del proyecto
    \item \textbf{Hito 4 (22/06)}: \acrshort{api} de control completado y funcionalidades avanzadas finalizadas
    \item \textbf{Hito 5 (23/06)}: Documentación técnica y tipado de código finalizados
\end{itemize}

\subsection{Adaptaciones en el seguimiento}
\label{subsec:adaptaciones_seguimiento}

Los cambios más significativos en el seguimiento reflejaron las limitaciones descubiertas durante la implementación. La migración temprana de Julia a Python ilustró la importancia de priorizar la productividad de desarrollo sobre el rendimiento teórico. Los dos fallos consecutivos de integración directa entre módulos revelaron limitaciones fundamentales en la aproximación arquitectónica inicial, llevando a la progresión desde integración monolítica con \acrshort{lsl} hacia refactorización \acrshort{solid} y finalmente hacia arquitectura \acrshort{api}.

\section{Análisis de desvíos y replanificación}
\label{sec:desvios_replanificacion}

El desarrollo experimentó desvíos significativos que requirieron adaptaciones metodológicas y replanificación de actividades.

\subsection{Principales desvíos identificados}
\label{subsec:principales_desvios}

Los desvíos más significativos reflejaron limitaciones en la estimación inicial de complejidad:

\begin{itemize}
    \item \textbf{Fallo de integración monolítica (17-19 abril)}: Reveló limitaciones arquitectónicas fundamentales no previstas en la planificación
    \item \textbf{Segundo fallo post-\acrshort{solid} (22-30 abril)}: Confirmó necesidad de cambio estratégico hacia arquitecturas \acrshort{api}
    \item \textbf{Pausa no planificada (1-13 mayo)}: Debida a compromisos académicos externos que afectaron la continuidad
    \item \textbf{Desarrollo de arquitectura \acrshort{api} (14-22 mayo)}: Solución emergente que requirió tiempo no contemplado inicialmente
    \item \textbf{Identificación problemas modo \textit{headless} (9-12 junio)}: Limitaciones detectadas tardíamente en el desarrollo
\end{itemize}

\subsection{Cronograma real ejecutado}
\label{subsec:cronograma_real}

El cronograma real muestra una evolución más compleja que la planificación inicial, con múltiples replanificaciones y desarrollo de funcionalidades no anticipadas.

\begin{sidewaysfigure}[p]
\centering
\resizebox{0.95\textwidth}{0.8\textheight}{
\begin{ganttchart}[
    canvas/.append style={fill=none, draw=black!5, line width=.75pt},
    hgrid style/.style={draw=black!5, line width=.75pt},
    vgrid={*1{draw=black!5, line width=.75pt}},
    title/.style={draw=none, fill=none},
    title label font=\bfseries\small,
    title label node/.append style={below=5pt},
    include title in canvas=false,
    bar label font=\mdseries\small\color{black!70},
    bar label node/.append style={left=2cm},
    bar/.append style={draw=none, fill=black!63},
    bar incomplete/.append style={fill=red!50},
    milestone/.append style={fill=yellow, draw=black, diamond},
    bar progress label font=\mdseries\small\color{black!70},
    group incomplete/.append style={fill=groupblue},
    group left shift=0,
    group right shift=0,
    group height=.5,
    group peaks tip position=0,
    group label node/.append style={left=1cm},
    group progress label font=\bfseries\small
  ]{1}{84}
  
  \gantttitle{Abril 2025}{21} \gantttitle{Mayo 2025}{21} \gantttitle{Junio 2025}{21} \gantttitle{Final}{21} \\
  \gantttitlelist{12,19,26}{7} \gantttitlelist{3,10,17}{7} \gantttitlelist{24,31,7}{7} \gantttitlelist{14,21,28}{7} \\
  
  \ganttgroup[progress=100]{Desarrollo módulos básicos}{1}{12} \\
  \ganttbar[progress=100]{Script \textit{inputs} + \acrshort{gui} PyQt5}{1}{7} \\
  \ganttbar[progress=100]{Módulo \acrshort{eeg} + \textit{brain map}}{8}{12} \\
  
  \ganttgroup[progress=100]{Primer intento integración}{13}{16} \\
  \ganttbar[progress=100, bar/.append style={fill=red!70}]{\acrshort{lsl} monolítico (fallido)}{13}{16} \\
  \ganttmilestone[milestone/.append style={fill=red}]{Fallo integración 1}{16} \\
  
  \ganttgroup[progress=100]{Refactorización \acrshort{solid}}{17}{19} \\
  \ganttbar[progress=100]{Reestructuración módulos}{17}{19} \\
  
  \ganttgroup[progress=100]{Segundo intento integración}{20}{28} \\
  \ganttbar[progress=100, bar/.append style={fill=red!70}]{Módulo integrador (fallido)}{20}{28} \\
  \ganttmilestone[milestone/.append style={fill=red}]{Fallo integración 2}{28} \\
  
  \ganttgroup[progress=0]{Pausa no planificada}{29}{42} \\
  \ganttbar[progress=0, bar/.append style={fill=gray!50}]{Compromisos académicos}{29}{42} \\
  
  \ganttgroup[progress=100]{Desarrollo \acrshort{api} y explorador web}{43}{56} \\
  \ganttbar[progress=100, bar/.append style={fill=green!70}]{\acrshort{api} acceso datos}{43}{49} \\
  \ganttbar[progress=100, bar/.append style={fill=green!70}]{Explorador web unificado}{50}{53} \\
  \ganttbar[progress=100, bar/.append style={fill=green!70}]{\acrshort{api} control \textit{inputs}}{54}{56} \\
  \ganttmilestone[milestone/.append style={fill=green}]{Arquitectura \acrshort{api} base exitosa}{56} \\
  
  \ganttgroup[progress=100]{\acrshort{cli} y optimizaciones}{57}{63} \\
  \ganttbar[progress=100]{\acrshort{cli} \textit{inputs} + \acrshort{ps4} + eventos}{57}{63} \\
  \ganttmilestone[milestone/.append style={fill=orange}]{Problemas \textit{headless} identificados}{61} \\
  
  \ganttgroup[progress=100]{Funcionalidades avanzadas \acrshort{eeg}}{64}{77} \\
  \ganttbar[progress=100]{\acrshort{api} control \acrshort{eeg} + \textit{MNE}}{64}{70} \\
  \ganttbar[progress=100]{\acrshort{ica} + \gls{wavelet} + filtros}{71}{77} \\
  \ganttmilestone[milestone/.append style={fill=green}]{\acrshort{api} control completo}{77} \\
  
  \ganttgroup[progress=50]{Integración final}{78}{80} \\
  \ganttbar[progress=50, bar/.append style={fill=orange!50}]{Módulo \acrshort{lsl} (pausado por \textit{bugs})}{78}{80} \\
  \ganttmilestone[milestone/.append style={fill=orange}]{Integración pausada}{80} \\
  
\end{ganttchart}
}
\caption{Diagrama de Gantt - Desarrollo real con desvíos}
\label{fig:gantt_real}
\end{sidewaysfigure}

La Figura~\ref{fig:gantt_real} revela una evolución significativamente más compleja que la planificación inicial, con múltiples intentos fallidos, replanificaciones y desarrollo de funcionalidades no anticipadas que finalmente resultaron en un sistema más robusto.

\subsection{Impacto económico de los desvíos}
\label{subsec:impacto_economico_desvios}

Los desvíos incrementaron el tiempo de desarrollo pero permitieron alcanzar un sistema más robusto y con mayor funcionalidad que la planificada inicialmente.

\begin{table}[htbp!]
  \centering
  \resizebox{\textwidth}{!}{
    \rowcolors{2}{white}{udcgray!25}
    \begin{tabular}{l|c|c|c}
    \rowcolor{udcpink!25}
    \textbf{Concepto} & \textbf{Planificado} & \textbf{Real} & \textbf{Desviación} \\
    \hline
    Duración del proyecto & 16 semanas & 18.5 semanas & +2.5 semanas \\
    Horas de desarrollo & 280 horas & 350 horas & +70 horas (+25\%) \\
    Coste de desarrollo & 3.360€ & 4.200€ & +840€ \\
    Alcance funcional & Básico planificado & Ampliado significativamente & +40\% funcionalidades \\
    \hline
    \end{tabular}
  }
  \caption{Comparación planificado vs. real del proyecto}
  \label{tab:comparacion_planificado_real}
\end{table}

La Tabla~\ref{tab:comparacion_planificado_real} presenta el análisis comparativo entre la planificación inicial y el desarrollo real, mostrando que aunque el proyecto requirió 25\% más tiempo del planificado, el alcance final superó significativamente los objetivos originales con un incremento del 40\% en funcionalidades.